\begin{IEEEbiography}[{\authorphoto{A1.jpg}}]{SONTAYA MANMAI}
(Student Member, IEEE) received the B.IET. degree in Industrial
Electrical Technology from King Mongkut's Institute of Technology
North Bangkok, Thailand, in 2000, and the M.EE. degree in electrical
power from Rajamangala University of Technology Thanyaburi (RMUTT),
Pathum Thani, Thailand, in 2012.

He is currently with the Department of Electrical Engineering,
Faculty of Engineering, Rajamangala University of Technology
Thanyaburi, Pathum Thani, Thailand. His research interests include
power and energy systems, renewable-energy integration, electric
vehicles, power-system stability assessment and prediction, and
smart and flexible power grids.
\end{IEEEbiography}


\begin{IEEEbiography}[{\authorphoto{A2.png}}]{PRAKASIT PRABPAL}
received the B.Eng. degree in Electrical Engineering from
Khon Kaen University (KKU), Khon Kaen, Thailand, in 2008,
the M.Eng. degree in Electrical and Computer Engineering from
Mahasarakham University (MSU), Maha Sarakham, Thailand, in 2011,
and the D.Eng. degree in Electrical Engineering from Rajamangala
University of Technology Thanyaburi (RMUTT), Pathum Thani,
Thailand, in 2021.

He is currently with the Department of Electrical Technology,
Udon Thani Technical College, Institute of Vocational Education
Northeastern Region 1, Udon Thani, Thailand. His research interests
include electrical power systems, power-system dynamics, energy
storage, electric vehicles, vehicle-to-grid applications, smart-grid
operation, and optimization techniques for modern power and energy
systems.
\end{IEEEbiography}


\begin{IEEEbiography}[{\authorphoto{A3.png}}]{YUTTANA KONGJEEN}
received the B.Eng. and M.Eng. degrees in Electrical Engineering
from Rajamangala University of Technology Isan (RMUTI),
Nakhon Ratchasima, Thailand, and the D.Eng. degree from
Rajamangala University of Technology Thanyaburi (RMUTT),
Pathum Thani, Thailand.

He is currently an Assistant Professor with the Department of
Electrical Engineering, Faculty of Engineering and Technology,
Rajamangala University of Technology Isan, Nakhon Ratchasima,
Thailand. He is also a member of the Intelligent Power System and
Energy Research (IPER) group. His research interests include
power-system optimization, electric vehicles, distribution systems,
energy storage, microgrids, energy management, power-system
stability, and renewable-energy integration.
\end{IEEEbiography}


\begin{IEEEbiography}[{\authorphoto{A4.jpg}}]{WASAN PHETPHIMOON}
received the M.Eng. and D.Eng. degrees in Electrical Engineering
from Rajamangala University of Technology Thanyaburi (RMUTT),
Pathum Thani, Thailand.

He is currently an Assistant Professor with the Department of
Electrical Engineering Technology, Faculty of Industrial Technology,
Kamphaeng Phet Rajabhat University, Kamphaeng Phet, Thailand.
His research interests include power electronics, high-frequency
power converters, renewable-energy conversion, photovoltaic systems,
battery-charging converters, bidirectional dc--dc converters,
electric vehicles, and energy-storage applications.
\end{IEEEbiography}


\begin{IEEEbiography}[{\authorphoto{A5.png}}]{KULSOMSUP YENCHAMCHALIT}
received the B.Ind.Tech. degree in Electrical Industrial Technology
from Thepsatri Rajabhat University, Thailand, in 2001, the M.Eng.
degree in Electrical Engineering from Mahanakorn University of
Technology, Thailand, in 2010, and the D.Eng. degree in Electrical
Engineering from Rajamangala University of Technology Thanyaburi,
Pathum Thani, Thailand, in 2022.

She is currently a Lecturer with the Department of Electrical
Industrial Technology, Faculty of Industrial Technology, Thepsatri
Rajabhat University, Lopburi, Thailand, where she also serves as
Deputy Dean for Administration. Her research interests include
electrical power systems, distributed energy resources, renewable-
energy integration, electric vehicles, and intelligent control
applications in modern power systems.
\end{IEEEbiography}


\begin{IEEEbiography}[{\authorphoto{A6.png}}]{SOMCHAT SONASANG}
(Member, IEEE) was born in Nongbua Lamphu, Thailand. He received
the B.Eng. degree in Electronics and Telecommunication Engineering
from Rajamangala University of Technology Isan (RMUTI),
Nakhon Ratchasima Campus, Thailand, in 2006, the M.Eng. degree
in Electrical Engineering from Khon Kaen University (KKU),
Khon Kaen, Thailand, in 2009, and the Ph.D. degree in Electrical
and Computer Engineering from Mahasarakham University (MSU),
Maha Sarakham, Thailand, in 2021.

He is currently an Associate Professor with the Faculty of
Industrial Technology, Nakhon Phanom University, Nakhon Phanom,
Thailand, and is affiliated with the Electrical Elements and
Wireless Devices Unit (EEWDU). His research interests include
microwave circuits and devices, electromagnetic material
characterization, nondestructive microwave sensing, antennas,
and applied electromagnetics. His recent research interests also
include renewable-energy systems, electric vehicles, vehicle-to-grid
integration, intelligent optimization, and computational methods
for modern power and energy systems.
\end{IEEEbiography}


\begin{IEEEbiography}[{\authorphoto{A7.jpg}}]{KRISCHONME BHUMKITTIPICH}
received the B.Eng. degree in Electrical Power Engineering
from Rajamangala University of Technology Thanyaburi (RMUTT),
Pathum Thani, Thailand, in 1997, the M.Eng. degree in
Electrical Engineering from Chulalongkorn University,
Bangkok, Thailand, in 2000, and the D.Eng. degree in Energy
from the Asian Institute of Technology (AIT), Pathum Thani,
Thailand, in 2008.

From 2002 to 2004, he was a Research Associate with the
Institute for Power Electronics and Electrical Drives (ISEA),
RWTH Aachen University, Aachen, Germany. He is currently
a Professor with the Department of Electrical Engineering,
Faculty of Engineering, Rajamangala University of Technology
Thanyaburi, Pathum Thani, Thailand, and serves as Vice President
for Academic and Research Affairs.

His research interests include electrical power systems,
power-system stability, renewable-energy integration, electric
vehicles, energy storage, smart grids, power electronics,
energy management, and optimization of modern power systems.
\end{IEEEbiography}
